% GENERATED FILE. Edit canonical lessons, then run npm run generate:books.
% canonical-source-hash: fnv1a64:3eedb892b8fa4f68
% canonical-lessons: RU-C506-gryadka, RU-C506-leyka, RU-C506-polot, RU-C506-sornyak, RU-C506-urozhay

\chapter{In the Vegetable Garden}
\label{ch:ru-in-the-vegetable-garden}
\hlchaptermodality{506}

\begin{chapteropening}
\textbf{By the end of this chapter:} \emph{I can talk about growing vegetables.}

Complete the last lesson of chapter 506: I can talk about growing vegetables.
\end{chapteropening}

\section[\texorpdfstring{\ru{грядка} (gryádka)}{грядка (gryádka)}]{\ru{грядка} (gryádka) — a garden bed}

\label{lesson:RU-C506-gryadka}



\begin{quote}
Before the new one: say the Russian for to bite, then the Russian for a catch.
\end{quote}

\subsection*{You'll want to know: \ru{грядка}}
\textbf{\ru{грядка}} (\emph{gryádka}) — \textquotedblleft{}a garden bed\textquotedblright{}.

\begin{quote}
\ru{На} \ru{этой} \ru{грядке} \ru{растёт} \ru{морковь}.
\end{quote}

\begin{quote}
\emph{Na étoy gryádke rastyót morkóv.}
\end{quote}

\begin{quote}
Carrots are growing in this bed.
\end{quote}

Feminine. A diminutive of \emph{\ru{гряда}} — a ridge.

\begin{cousinweb}
\emph{\ru{грядка}} is a diminutive of \emph{\ru{гряда}} 'a ridge, a row'; EtymDB has no entry.
\end{cousinweb}

\begin{grammarlens}[title={What you've built}]
One more word for telling, explaining and describing at B1.
\end{grammarlens}

\subsection*{Your turn}
\begin{itemize}
\raggedright
  \item \emph{Say it:} \emph{gryádka}
  \item \emph{Say it:} \emph{gryádka}, once more
  \item \emph{Say it:} a sentence of your own with \emph{gryádka}
  \item \emph{Recall:} say the Russian for to bite, then the Russian for a catch, then say \emph{gryádka} again
\end{itemize}

\begin{tcolorbox}[breakable,colback=teal!4,colframe=teal!35!black,title={Before you move on}]
What does \ru{грядка} mean? (\textquotedblleft{}A garden bed\textquotedblright{}.) Say it once more.
\end{tcolorbox}
\section[\texorpdfstring{\ru{лейка} (léyka)}{лейка (léyka)}]{\ru{лейка} (léyka) — a watering can}

\label{lesson:RU-C506-leyka}



\begin{quote}
Before the new one: say the Russian for a catch, then the Russian for a garden bed.
\end{quote}

\subsection*{You'll want to know: \ru{лейка}}
\textbf{\ru{лейка}} (\emph{léyka}) — \textquotedblleft{}a watering can\textquotedblright{}.

\begin{quote}
\ru{Девочка} \ru{поливает} \ru{цветы} \ru{из} \ru{лейки}.
\end{quote}

\begin{quote}
\emph{Dévochka poliváyet tsvetý iz léyki.}
\end{quote}

\begin{quote}
The girl is watering the flowers with a watering can.
\end{quote}

Feminine.

\begin{cousinweb}
\emph{\ru{лейка}} is built on the root of \emph{\ru{лить}} 'to pour'; EtymDB has no entry.
\end{cousinweb}

\begin{grammarlens}[title={What you've built}]
One more word for telling, explaining and describing at B1.
\end{grammarlens}

\subsection*{Your turn}
\begin{itemize}
\raggedright
  \item \emph{Say it:} \emph{léyka}
  \item \emph{Say it:} \emph{léyka}, once more
  \item \emph{Say it:} a sentence of your own with \emph{léyka}
  \item \emph{Recall:} say the Russian for a catch, then the Russian for a garden bed, then say \emph{léyka} again
\end{itemize}

\begin{tcolorbox}[breakable,colback=teal!4,colframe=teal!35!black,title={Before you move on}]
What does \ru{лейка} mean? (\textquotedblleft{}A watering can\textquotedblright{}.) Say it once more.
\end{tcolorbox}
\section[\texorpdfstring{\ru{полоть} (polót)}{полоть (polót)}]{\ru{полоть} (polót) — to weed}

\label{lesson:RU-C506-polot}



\begin{quote}
Before the new one: say the Russian for a garden bed, then the Russian for a watering can.
\end{quote}

\subsection*{You'll want to know: \ru{полоть}}
\textbf{\ru{полоть}} (\emph{polót}) — \textquotedblleft{}to weed\textquotedblright{}.

\begin{quote}
\ru{Летом} \ru{бабушка} \ru{часто} \ru{полет} \ru{огород}.
\end{quote}

\begin{quote}
\emph{Létom bábushka chásto pólet ogoród.}
\end{quote}

\begin{quote}
In the summer Grandma often weeds the vegetable garden.
\end{quote}

Imperfective; present \ru{полю}, \ru{полешь}. Perfective \ru{прополоть}.

\begin{cousinweb}
EtymDB lists \emph{\ru{полоть}} without a source; it is usually traced to Proto-Slavic \emph{pelti}.
\end{cousinweb}

\begin{grammarlens}[title={What you've built}]
One more word for telling, explaining and describing at B1.
\end{grammarlens}

\subsection*{Your turn}
\begin{itemize}
\raggedright
  \item \emph{Say it:} \emph{polót}
  \item \emph{Say it:} \emph{polót}, once more
  \item \emph{Say it:} a sentence of your own with \emph{polót}
  \item \emph{Recall:} say the Russian for a garden bed, then the Russian for a watering can, then say \emph{polót} again
\end{itemize}

\begin{tcolorbox}[breakable,colback=teal!4,colframe=teal!35!black,title={Before you move on}]
What does \ru{полоть} mean? (\textquotedblleft{}To weed\textquotedblright{}.) Say it once more.
\end{tcolorbox}
\section[\texorpdfstring{\ru{сорняк} (sornyák)}{сорняк (sornyák)}]{\ru{сорняк} (sornyák) — a weed}

\label{lesson:RU-C506-sornyak}



\begin{quote}
Before the new one: say the Russian for a watering can, then the Russian for to weed.
\end{quote}

\subsection*{You'll want to know: \ru{сорняк}}
\textbf{\ru{сорняк}} (\emph{sornyák}) — \textquotedblleft{}a weed\textquotedblright{}.

\begin{quote}
\ru{После} \ru{дождя} \ru{сорняки} \ru{растут} \ru{очень} \ru{быстро}.
\end{quote}

\begin{quote}
\emph{Pósle dozhdyá sornyakí rastút óchen býstro.}
\end{quote}

\begin{quote}
After rain the weeds grow very fast.
\end{quote}

Masculine; genitive \ru{сорняка}.

\begin{cousinweb}
\emph{\ru{сорняк}} is built on \emph{\ru{сор}} 'litter, rubbish' with \emph{-\ru{няк}}; EtymDB has no entry.
\end{cousinweb}

\begin{grammarlens}[title={What you've built}]
One more word for telling, explaining and describing at B1.
\end{grammarlens}

\subsection*{Your turn}
\begin{itemize}
\raggedright
  \item \emph{Say it:} \emph{sornyák}
  \item \emph{Say it:} \emph{sornyák}, once more
  \item \emph{Say it:} a sentence of your own with \emph{sornyák}
  \item \emph{Recall:} say the Russian for a watering can, then the Russian for to weed, then say \emph{sornyák} again
\end{itemize}

\begin{tcolorbox}[breakable,colback=teal!4,colframe=teal!35!black,title={Before you move on}]
What does \ru{сорняк} mean? (\textquotedblleft{}A weed\textquotedblright{}.) Say it once more.
\end{tcolorbox}
\section[\texorpdfstring{\ru{урожай} (urozháy)}{урожай (urozháy)}]{\ru{урожай} (urozháy) — a harvest, a crop}

\label{lesson:RU-C506-urozhay}



\begin{quote}
Before the new one: say the Russian for to weed, then the Russian for a weed.
\end{quote}

\subsection*{You'll want to know: \ru{урожай}}
\textbf{\ru{урожай}} (\emph{urozháy}) — \textquotedblleft{}a harvest, a crop\textquotedblright{}.

\begin{quote}
\ru{В} \ru{этом} \ru{году} \ru{у} \ru{нас} \ru{хороший} \ru{урожай} \ru{яблок}.
\end{quote}

\begin{quote}
\emph{V étom godú u nas khoróshiy urozháy yáblok.}
\end{quote}

\begin{quote}
We have a good apple crop this year.
\end{quote}

Masculine.

\begin{cousinweb}
\emph{\ru{урожай}} is inherited from Proto-Slavic \emph{urodjaj\ru{ь}} (EtymDB), built on the root of \emph{\ru{родить}} 'to bear, to give birth'.
\end{cousinweb}

\begin{grammarlens}[title={What you've built}]
One more word for telling, explaining and describing at B1.
\end{grammarlens}

\subsection*{Your turn}
\begin{itemize}
\raggedright
  \item \emph{Say it:} \emph{urozháy}
  \item \emph{Say it:} \emph{urozháy}, once more
  \item \emph{Say it:} a sentence of your own with \emph{urozháy}
  \item \emph{Recall:} say the Russian for to weed, then the Russian for a weed, then say \emph{urozháy} again
\end{itemize}

\begin{tcolorbox}[breakable,colback=teal!4,colframe=teal!35!black,title={Before you move on}]
What does \ru{урожай} mean? (\textquotedblleft{}A harvest, a crop\textquotedblright{}.) Say it once more.
\end{tcolorbox}
